\documentclass[11pt]{article}
\usepackage[margin=1in]{geometry}
\usepackage{amsmath,amssymb,amsthm}
\usepackage[utf8]{inputenc}
\usepackage{listings}
\usepackage{hyperref}
\usepackage{xcolor}
\usepackage{textcomp}

\definecolor{leanblue}{rgb}{0.2,0.4,0.7}

\lstdefinelanguage{Lean}{
  keywords={theorem,lemma,def,axiom,structure,class,instance,where,by,sorry,simp,exact,intro,apply,have,let,fun,match,if,then,else,import,open,namespace,end,section,variable,inductive,noncomputable,abbrev,deriving},
  keywordstyle=\color{leanblue}\bfseries,
  comment=[l]{--},
  morecomment=[s]{/-}{-/},
  string=[b]",
  basicstyle=\ttfamily\small,
  breaklines=true,
  extendedchars=true,
  inputencoding=utf8,
  literate=
    {≥}{{$\geq$}}1
    {≤}{{$\leq$}}1
    {→}{{$\to$}}1
    {←}{{$\leftarrow$}}1
    {∧}{{$\wedge$}}1
    {∨}{{$\vee$}}1
    {¬}{{$\neg$}}1
    {∀}{{$\forall$}}1
    {∃}{{$\exists$}}1
    {⊆}{{$\subseteq$}}1
    {⊂}{{$\subset$}}1
    {∈}{{$\in$}}1
    {∉}{{$\notin$}}1
    {×}{{$\times$}}1
    {⊗}{{$\otimes$}}1
    {▸}{{$\blacktriangleright$}}1
    {⟨}{{$\langle$}}1
    {⟩}{{$\rangle$}}1
    {λ}{{$\lambda$}}1
    {α}{{$\alpha$}}1
    {β}{{$\beta$}}1
    {σ}{{$\sigma$}}1
    {θ}{{$\theta$}}1
    {ε}{{$\varepsilon$}}1
    {μ}{{$\mu$}}1
    {π}{{$\pi$}}1
    {Σ}{{$\Sigma$}}1
    {Π}{{$\Pi$}}1
    {▶}{{$\triangleright$}}1
    {⊔}{{$\sqcup$}}1
    {⊓}{{$\sqcap$}}1
    {⊥}{{$\bot$}}1
    {⊤}{{$\top$}}1
    {∅}{{$\emptyset$}}1
    {∪}{{$\cup$}}1
    {∩}{{$\cap$}}1
    {≠}{{$\neq$}}1
    {ℕ}{{$\mathbb{N}$}}1
    {₀}{{$_0$}}1
    {₁}{{$_1$}}1
    {₂}{{$_2$}}1
    {|>}{{$|{\!>}$}}2
    {·}{{$\cdot$}}1
    {—}{{---}}1
    {∘}{{$\circ$}}1
    {√}{{$\sqrt{}$}}1
    {∝}{{$\propto$}}1
    {≈}{{$\approx$}}1
    {═}{{=}}1
    {⟹}{{$\Longrightarrow$}}1,
}

\newtheorem{theorem}{Theorem}
\newtheorem{lemma}[theorem]{Lemma}
\newtheorem{definition}[theorem]{Definition}
\newtheorem{axiom_stmt}[theorem]{Axiom}

\title{Formal Verification: Threshold Protocol Composition}
\author{Zach Kelling, Lux Industries}
\date{December 2025}

\begin{document}
\maketitle

\begin{abstract}
We prove that the Lux threshold protocols compose safely: FROST, CMP (CGGMP21), TFHE, Ringtail, and BLS all share the same $t$-of-$n$ access structure via LSS. Key material from one protocol does not leak to another, and all protocols can be initialized in a single coordinated ceremony.
\end{abstract}

\section{Introduction}
The composition theorem ensures that the Lux threshold protocol stack is coherent. A single key ceremony can initialize shares for all protocols simultaneously, and compromising one protocol's key material does not affect the others. The dependency structure is a flat fan-out from LSS.

\section{Definitions}
\begin{definition}[Protocol]
Protocol identifier: FROST, CMP, TFHE, Ringtail, BLS, LSS.
\end{definition}

\begin{definition}[UnifiedParams]
Shared threshold parameters across all protocols.
\end{definition}

\begin{definition}[ProtocolInstance]
A protocol bound to unified parameters.
\end{definition}

\begin{definition}[KeyMaterial]
Protocol-specific key material (no cross-contamination).
\end{definition}


\section{Theorems and Axioms}
\begin{theorem}[\texttt{uniform\_access\_structure}]
All protocols instantiated with the same \texttt{UnifiedParams} share the same access structure.
\end{theorem}

\begin{theorem}[\texttt{threshold\_consistent}]
The threshold $t$ is identical across FROST, CMP, and TFHE when using unified parameters.
\end{theorem}

\begin{theorem}[\texttt{six\_protocols}]
The Lux stack includes exactly 6 threshold protocols.
\end{theorem}

\begin{theorem}[\texttt{multi\_keygen\_uniform}]
All instances from \texttt{multiProtocolKeygen} share the same access structure.
\end{theorem}

\begin{theorem}[\texttt{lss\_no\_deps}]
LSS has no protocol dependencies (it is the base layer).
\end{theorem}

\begin{theorem}[\texttt{all\_depend\_on\_lss}]
Every non-LSS protocol depends on LSS.
\end{theorem}
\begin{proof}
By case analysis on all protocol constructors.
\end{proof}


\section{Lean Source}
\lstinputlisting[language=Lean]{../../formal/lean/Crypto/Threshold/Composition.lean}

\section{References}
\noindent Source code: \url{https://github.com/luxfi/proofs/blob/main/lean/Crypto/Threshold/Composition.lean}\\
\noindent White papers: \url{https://papers.lux.network}

\noindent Related white papers:
\begin{itemize}
  \item \texttt{lux-universal-threshold-signatures.tex}
\end{itemize}

\end{document}
